% GENERATED FILE. Edit canonical lessons, then run npm run generate:books.
% canonical-source-hash: fnv1a64:d9be175c1aa64d7f
% canonical-lessons: BN-C467-thot, BN-C467-choyal, BN-C467-mari, BN-C467-kontho, BN-C467-kochkano

\chapter{Face and Mouth}
\label{ch:bn-face-and-mouth}
\hlchaptermodality{467}

\begin{chapteropening}
\textbf{By the end of this chapter:} \emph{I can name the lips, jaw, gums and voice, and say that someone frowned.}

Complete the last lesson of chapter 467: I can name the lips, jaw, gums and voice, and say that someone frowned.
\end{chapteropening}

\section[\texorpdfstring{\bn{ঠোঁট} (ṭhõṭ)}{ঠোঁট (ṭhõṭ)}]{\bn{ঠোঁট} (ṭhõṭ) — a lip}

\label{lesson:BN-C467-thot}



\begin{quote}
Before the new one: say the Bengali for the back, then the Bengali for the nape.
\end{quote}

\subsection*{You'll want to know: \bn{ঠোঁট}}
\textbf{\bn{ঠোঁট}} — \emph{ṭhõṭ} — \textquotedblleft{}a lip\textquotedblright{}.

\begin{quote}
\bn{শীতে} \bn{মীরার} \bn{ঠোঁট} \bn{শুকনো} \bn{হয়ে} \bn{যায়।}
\end{quote}

\begin{quote}
\emph{shīte mīrār ṭhõṭ shukno hôye jāy.}
\end{quote}

\begin{quote}
In winter Mira's lips get dry.
\end{quote}

\begin{cousinweb}
\textbf{\bn{ঠোঁট}} (\emph{ṭhõṭ}) is listed by CDIAL for Bengali, as a bird's beak, under Sanskrit \emph{tuṇḍa}, \textquotedblleft{}beak, snout\textquotedblright{} (CDIAL 5853); everyday Bengali uses it for the human lip.
\end{cousinweb}

\begin{grammarlens}[title={What you've built}]
One more word for describing what you lived and felt.
\end{grammarlens}

\subsection*{Your turn}
\begin{itemize}
\raggedright
  \item \emph{Say it:} \emph{ṭhõṭ}
  \item \emph{Say it:} \emph{ṭhõṭ}, once more
  \item \emph{Say it:} \emph{shīte mīrār ṭhõṭ shukno hôye jāy.} — In winter Mira's lips get dry.
  \item \emph{Recall:} say the Bengali for the back, then the Bengali for the nape, then say \emph{ṭhõṭ} again
\end{itemize}

\begin{tcolorbox}[breakable,colback=teal!4,colframe=teal!35!black,title={Before you move on}]
What does \bn{ঠোঁট} mean? (\textquotedblleft{}A lip\textquotedblright{}.) Say it once more.
\end{tcolorbox}
\section[\texorpdfstring{\bn{চোয়াল} (choyāl)}{চোয়াল (choyāl)}]{\bn{চোয়াল} (choyāl) — the jaw}

\label{lesson:BN-C467-choyal}



\begin{quote}
Before the new one: say the Bengali for the nape, then the Bengali for a lip.
\end{quote}

\subsection*{You'll want to know: \bn{চোয়াল}}
\textbf{\bn{চোয়াল}} — \emph{choyāl} — \textquotedblleft{}the jaw\textquotedblright{}.

\begin{quote}
\bn{শক্ত} \bn{মাংস} \bn{চিবিয়ে} \bn{আমার} \bn{চোয়াল} \bn{ব্যথা} \bn{করছে।}
\end{quote}

\begin{quote}
\emph{shôktô māngsô chibiye āmār choyāl bæthā kôrchhe.}
\end{quote}

\begin{quote}
Chewing the tough meat has made my jaw ache.
\end{quote}

\begin{cousinweb}
\textbf{\bn{চোয়াল}} (\emph{choyāl}) is a Bengali word for the jaw whose origin is not settled.
\end{cousinweb}

\begin{grammarlens}[title={What you've built}]
One more word for describing what you lived and felt.
\end{grammarlens}

\subsection*{Your turn}
\begin{itemize}
\raggedright
  \item \emph{Say it:} \emph{choyāl}
  \item \emph{Say it:} \emph{choyāl}, once more
  \item \emph{Say it:} \emph{shôktô māngsô chibiye āmār choyāl bæthā kôrchhe.} — Chewing the tough meat has made my jaw ache.
  \item \emph{Recall:} say the Bengali for the nape, then the Bengali for a lip, then say \emph{choyāl} again
\end{itemize}

\begin{tcolorbox}[breakable,colback=teal!4,colframe=teal!35!black,title={Before you move on}]
What does \bn{চোয়াল} mean? (\textquotedblleft{}The jaw\textquotedblright{}.) Say it once more.
\end{tcolorbox}
\section[\texorpdfstring{\bn{মাড়ি} (māṛi)}{মাড়ি (māṛi)}]{\bn{মাড়ি} (māṛi) — the gums}

\label{lesson:BN-C467-mari}



\begin{quote}
Before the new one: say the Bengali for a lip, then the Bengali for the jaw.
\end{quote}

\subsection*{You'll want to know: \bn{মাড়ি}}
\textbf{\bn{মাড়ি}} — \emph{māṛi} — \textquotedblleft{}the gums\textquotedblright{}.

\begin{quote}
\bn{অরুণের} \bn{মাড়ি} \bn{থেকে} \bn{রক্ত} \bn{পড়ছে।}
\end{quote}

\begin{quote}
\emph{ôruner māṛi theke rôkto pôṛchhe.}
\end{quote}

\begin{quote}
Arun's gums are bleeding.
\end{quote}

\begin{cousinweb}
\textbf{\bn{মাড়ি}} (\emph{māṛi}) is listed by CDIAL for Bengali under a reconstructed \emph{māṁsapuṭaka}, \textquotedblleft{}the gums\textquotedblright{} (CDIAL 9984), a compound of the words for flesh and for a fold.
\end{cousinweb}

\begin{grammarlens}[title={What you've built}]
One more word for describing what you lived and felt.
\end{grammarlens}

\subsection*{Your turn}
\begin{itemize}
\raggedright
  \item \emph{Say it:} \emph{māṛi}
  \item \emph{Say it:} \emph{māṛi}, once more
  \item \emph{Say it:} \emph{ôruner māṛi theke rôkto pôṛchhe.} — Arun's gums are bleeding.
  \item \emph{Recall:} say the Bengali for a lip, then the Bengali for the jaw, then say \emph{māṛi} again
\end{itemize}

\begin{tcolorbox}[breakable,colback=teal!4,colframe=teal!35!black,title={Before you move on}]
What does \bn{মাড়ি} mean? (\textquotedblleft{}The gums\textquotedblright{}.) Say it once more.
\end{tcolorbox}
\section[\texorpdfstring{\bn{কণ্ঠ} (kônṭhô)}{কণ্ঠ (kônṭhô)}]{\bn{কণ্ঠ} (kônṭhô) — the voice; the throat}

\label{lesson:BN-C467-kontho}



\begin{quote}
Before the new one: say the Bengali for the jaw, then the Bengali for the gums.
\end{quote}

\subsection*{You'll want to know: \bn{কণ্ঠ}}
\textbf{\bn{কণ্ঠ}} — \emph{kônṭhô} — \textquotedblleft{}the voice; the throat\textquotedblright{}.

\begin{quote}
\bn{গায়িকার} \bn{কণ্ঠ} \bn{সুন্দর।}
\end{quote}

\begin{quote}
\emph{gāyikār kônṭhô sundôr.}
\end{quote}

\begin{quote}
The singer has a beautiful voice.
\end{quote}

\begin{cousinweb}
\textbf{\bn{কণ্ঠ}} (\emph{kônṭhô}) is a tatsama, Sanskrit \emph{kaṇṭha}, \textquotedblleft{}the throat, the neck\textquotedblright{} and \textquotedblleft{}the voice\textquotedblright{} (Monier-Williams).
\end{cousinweb}

\begin{grammarlens}[title={What you've built}]
One more word for describing what you lived and felt.
\end{grammarlens}

\subsection*{Your turn}
\begin{itemize}
\raggedright
  \item \emph{Say it:} \emph{kônṭhô}
  \item \emph{Say it:} \emph{kônṭhô}, once more
  \item \emph{Say it:} \emph{gāyikār kônṭhô sundôr.} — The singer has a beautiful voice.
  \item \emph{Recall:} say the Bengali for the jaw, then the Bengali for the gums, then say \emph{kônṭhô} again
\end{itemize}

\begin{tcolorbox}[breakable,colback=teal!4,colframe=teal!35!black,title={Before you move on}]
What does \bn{কণ্ঠ} mean? (\textquotedblleft{}The voice; the throat\textquotedblright{}.) Say it once more.
\end{tcolorbox}
\section[\texorpdfstring{\bn{কোঁচকানো} (kõchkāno)}{কোঁচকানো (kõchkāno)}]{\bn{কোঁচকানো} (kõchkāno) — to wrinkle; to knit one's brows}

\label{lesson:BN-C467-kochkano}



\begin{quote}
Before the new one: say the Bengali for the gums, then the Bengali for the voice.
\end{quote}

\subsection*{You'll want to know: \bn{কোঁচকানো}}
\textbf{\bn{কোঁচকানো}} — \emph{kõchkāno} — \textquotedblleft{}to wrinkle; to knit one's brows\textquotedblright{}.

\begin{quote}
\bn{খারাপ} \bn{খবরে} \bn{বাবা} \bn{ভুরু} \bn{কোঁচকালেন।}
\end{quote}

\begin{quote}
\emph{khārāp khôbôre bābā bhuru kõchkālen.}
\end{quote}

\begin{quote}
Father knitted his brows at the bad news.
\end{quote}

Past (he/she …): \bn{তিনি} \bn{কপাল} \bn{কোঁচকালেন}, \emph{tini kopāl kõchkālen}.

\begin{cousinweb}
\textbf{\bn{কোঁচকানো}} (\emph{kõchkāno}) is a Bengali verb of puckering and wrinkling whose origin is not settled.
\end{cousinweb}

\begin{grammarlens}[title={What you've built}]
One more word for describing what you lived and felt.
\end{grammarlens}

\subsection*{Your turn}
\begin{itemize}
\raggedright
  \item \emph{Say it:} \emph{kõchkāno}
  \item \emph{Say it:} \emph{kõchkāno}, once more
  \item \emph{Say it:} \emph{khārāp khôbôre bābā bhuru kõchkālen.} — Father knitted his brows at the bad news.
  \item \emph{Recall:} say the Bengali for the gums, then the Bengali for the voice, then say \emph{kõchkāno} again
\end{itemize}

\begin{tcolorbox}[breakable,colback=teal!4,colframe=teal!35!black,title={Before you move on}]
What does \bn{কোঁচকানো} mean? (\textquotedblleft{}To wrinkle; to knit one's brows\textquotedblright{}.) Say it once more.
\end{tcolorbox}
